\documentclass[10pt,a4paper]{amsart}
\usepackage[margin=19mm]{geometry}
\usepackage{amsmath,amssymb,mathtools}
\usepackage[scr=rsfs,cal=euler]{mathalfa}
\usepackage{xfrac}
\usepackage[unicode,hidelinks,bookmarks=false]{hyperref}
\newcommand{\ie}{i.e.\ }
\newcommand{\cf}{cf.\ }
\newcommand{\resp}{respectively\ }
\newcommand{\Subsection}{\subsection}
\providecommand{\nobreakdash}{\nobreak\hbox{-}}
\setlength{\emergencystretch}{2em}
\multlinegap0pt
\usepackage{mathtools}
\usepackage{amssymb}
\usepackage{bm}
\usepackage{dsfont}
\newtheorem{proposition}{Proposition}
\newtheorem{lemma}{Lemma}
\newcommand{\G}{\mathbb{G}}
\renewcommand{\P}{\mathbb{P}}
\newcommand{\W}{\mathcal{W}}
\newcommand{\cbar}{\overline{C}}
\renewcommand{\geq}{\geqslant}
\renewcommand{\leq}{\leqslant}
\def\mathbi#1{\textbf{\em #1}}
\newcommand{\tu}[1]{\textup{#1}}
\title{Upper tail bounds for irregular graphs}
\author{Anirban Basak, Shaibal Karmakar}
\date{Source: arXiv:2503.05311v4 (CC BY 4.0)}
\begin{document}
\maketitle

\noindent\emph{Source and licence.} Upper tail bounds for irregular graphs. Version arXiv:2503.05311v4 (CC BY 4.0), distributed by arXiv under the Creative Commons Attribution 4.0 International licence, http://creativecommons.org/licenses/by/4.0/, as recorded in the arXiv licence metadata for this exact version and shown on the abstract page https://arxiv.org/abs/2503.05311v4. Full licence notice: \texttt{licenses/arxiv-2503.05311-CC-BY-4.0.txt}.\par\smallskip

\noindent\textit{Editorial scope.} Counted unit: the article's proposition proposition (source0.tex lines 1133--1144). The article's complete original proof of it is delivered with it (source0.tex lines 1365--1435, one \texttt{proof} environment of 71 lines, which the article places away from the statement). No mathematics is written, repaired or re-derived: the statement and the proof are the article's own lines, in the article's own notation, and no retained line is substituted or re-flowed.  Retained uncounted as the source-local context the proof needs: lines 1212--1215; lines 1246--1258; lines 1274--1281; lines 1326--1332.  Nothing was omitted from any retained span, so the package carries no omission record.  No retained line contains a citation, so the wrapper carries no bibliography and no bibliography entry.  No retained line contains a figure environment, an image-inclusion command or drawing code, so no image is delivered and the package declares no asset dependency.  Editorial formatting only, in addition: an AMS \texttt{amsart} preamble that restates the article's own declarations and macros used by the retained lines, namely the environments \texttt{proposition}, \texttt{lemma} on separate counters, together with the standard packages those lines need.  The retained environments print in this document as Proposition 1, Lemma 1 in order of appearance, whereas the article prints the same items with the numbering of its own full document.  The complete native source of the article as distributed in the arXiv e-print for this exact version is bundled unchanged as \texttt{sources/174376/source0.tex}.
			\begin{equation}
				\label{bipartite argument}
				\text{deg}_{\tu{G}}(v)\geq \varepsilon\cdot \text{deg}_{\tu{G}}(u)\cdot \text{deg}_{\tu{G}}(v)\geq \varepsilon \frac{\widetilde{c}_{0}n^{1+1/r}p}{(\log n)^{1/(r-1)}} \geq \frac{2}{\varepsilon},
			\end{equation}

			\begin{lemma}
				\label{counting cores}
				Fix $\varepsilon\in (0,1/2)$ and $r\geq 2$. %For any core graph define 
				%    \begin{equation*}
					%    \W(\tu{G}) \coloneqq \left\{v\in V(\tu{G}): \textup{deg}_{\tu{G}}(v)\leq \frac{1}{\varepsilon}\right\}.
					%    \end{equation*}
				%    Set $\overline{\W(\tu{G})} \coloneqq V(\tu{G})\setminus \W$. 
				%Further 
				Let $\mathcal{N}(\mathbi{e}, \textbf{w}, \varepsilon)$ be the number of core graphs with $e(\tu{G})=\mathbi{e}$ and $|\W(\tu{G})|=\textbf{w}$. Then for any $p\in (0,1)$ satisfying $n^{-1-1/r}(\log n)^{1/(r-1)}\ll p\lesssim n^{-1/r}$ and large enough $n$,
				\begin{equation*}
					\mathcal{N}(\mathbi{e}, \textbf{w}, \varepsilon)\leq {n\choose \textbf{w}}\cdot \exp\left(17r\varepsilon \mathbi{e}\log(1/p)\right).
				\end{equation*}
			\end{lemma}

			\begin{lemma}
				\label{global embedding bound}
				Assume $t\geq 2$.
				For every graph $\tu{G}\subseteq K_n$, 
				\begin{equation*}
					N(K_{1,t}, \tu{G})\leq \left(e(\tu{G})\right)^{t}.
				\end{equation*}
			\end{lemma}

		\begin{lemma}
			\label{s2: improved geb }
			Let $\tu{G}$ be a core graph on $n$ vertices. Assume $np^{r}\rightarrow \rho\in (0,\infty)$. Then,
			\begin{equation*}
				e(\tu{G})\geq (1-\varepsilon)\left(\lfloor\delta\rho(1-5\varepsilon)\rfloor + \{\delta\rho(1-5\varepsilon)\}^{1/r}\right) \cdot n.
			\end{equation*}
		\end{lemma}

			\begin{proposition}
				\label{main proposition}
				Assume $(\log n)^{1/(r-1)}\ll n^{1+1/r}p$. If $np^{r}\ll 1$, then for large enough $n$, we have that
				\begin{equation*}
					\P\left(\G(n,p)\textup{ contains a core graph}\right)\leq \exp\left(-(1-\mathfrak{f}(\varepsilon))\frac{\delta^{1/r}}{r}n^{1+1/r}p\log n\right),
				\end{equation*}
				and if $np^{r}\rightarrow \rho\in (0,\infty)$, then for large enough $n$, we have
				\begin{equation}\label{eq:p-sim-bdry}
					\P\left(\G(n,p)\textup{ contains a core graph}\right) \leq \exp\left(-(1-\mathfrak{f}(\varepsilon))\left(\lfloor \delta \rho\rfloor+\left\{\delta \rho\right\}^{1/r}\right)\frac{1}{r}n\log n\right),
				\end{equation}
				for some nonnegative function $\mathfrak{f}(\cdot)$ such that $\lim_{\varepsilon\downarrow 0}\mathfrak{f}(\varepsilon)=0$.
			\end{proposition}

		\begin{proof}[Proof of Proposition \ref{main proposition}]
			By Lemma \ref{global embedding bound} for any core graph $\tu{G}$ %, number of edges $e(\tu{G})$ must satisfy 
			\begin{equation*}
				e_{\min} \coloneqq \delta^{1/r}(1-3\varepsilon)^{1/r}n^{1+1/r}p \leq e(\tu{G})\leq \cbar n^{1+1/r}p\log(1/p) \eqqcolon e_{\max}.
			\end{equation*}
			%where the lower bound comes from fact that a core graph must contain at least $\delta(1-3\varepsilon)n^{r+1}p^r$ copies of $K_{1,r}$ and then using Lemma \ref{global embedding bound}. Set $e_{min} \coloneqq (1-3\varepsilon)^{1/r}\delta^{1/r}n^{1+1/r}p$ and $e_{max} \coloneqq \cbar n^{1+1/r}p\log(1/p)$. 
			%Also note from the discussion in \eqref{bipartite argument} 
			As $\tu{G}_{\mathcal{W}}$ is bipartite we also note that the set $\W(\tu{G})$ can have at most $e(\tu{G})$ elements.
			Fix any $e_{\min}\leq \mathbi{e} \leq e_{\max}$ and $\mathbi{w}\leq \mathbi{e}$. We claim that %Then by using Lemma \ref{counting cores} we get
			\begin{multline}
				\label{main proof eq 1}
				\P\left(\exists \tu{G}\subset \G(n,p):\ e(\tu{G})=\mathbi{e}\text{ and } |\W(\tu{G})|=\mathbi{w}\right)\leq {n\choose \mathbi{w}}\cdot p^{-17r\varepsilon \mathbi{e}}\cdot p^{\mathbi{e}} \\
				\leq \exp\left(-(1-36r^{2}\varepsilon)(1-3\varepsilon)^{1/r}\frac{\delta^{1/r}}{r}n^{1+1/r}p\log n\right).
			\end{multline}
			The first inequality is immediate from Lemma \ref{counting cores}. To prove the second inequality we split the range of  ${\bm w}$ and ${\bm e}$ into two cases: $\mathbi{w}\leq \varepsilon \mathbi{e}$ and $\mathbi{w}\geq \varepsilon \mathbi{e} $. In the first case, we have
			\begin{equation}
				\label{main proof eq 2}
				{n\choose \mathbi{w}}\cdot p^{(1-17r\varepsilon)\mathbi{e}}\leq n^{\mathbi{w}}\cdot p^{(1-17r\varepsilon)\mathbi{e}}\leq n^{\varepsilon \mathbi{e}}\cdot p^{(1-17r\varepsilon)\mathbi{e}} \leq n^{-(1-19r\varepsilon)\mathbi{e}/r},
			\end{equation}
			where in the last step we used $p\lesssim  n^{-1/r}$.
			In the second case, using the lower bound $\mathbi{w}\geq \varepsilon \mathbi{e}\geq \varepsilon e_{\min}$ we obtain that, for all large $n$
			\begin{equation*}
				{n\choose \mathbi{w}} \cdot p^{(1-17r\varepsilon)\mathbi{e}} \leq \left(\frac{en}{\mathbi{w}}\right)^{\mathbi{w}} \cdot p^{(1-17r\varepsilon)\mathbi{e}}\leq \left(\frac{e}{\delta^{1/r}(1-3\varepsilon)^{1/r}}\right)^{\mathbi{w}}\cdot n^{-\mathbi{w}/r}\cdot p^{(1-17r\varepsilon)\mathbi{e}-\mathbi{w}}\leq n^{-\mathbi{w}/r} \cdot p^{(1-18r\varepsilon)\mathbi{e}-\mathbi{w}},
			\end{equation*}
			where in the last step we have used $\mathbi{w}\leq \mathbi{e}$ and $p\ll1$. 
			%Therefore, we have for large $n$
			%\begin{equation*}
			%    {n\choose \mathbi{w}}\cdot p^{(1-7r\varepsilon)\mathbi{e}}\leq n^{-\mathbi{w}/r}\cdot p^{(1-18r\varepsilon)\mathbi{e}-\mathbi{w}}.
			%\end{equation*}
			To complete the proof of \eqref{main proof eq 1} we further split into two sub cases: $\mathbi{w}\leq (1-18r\varepsilon)\mathbi{e}$ and $\mathbi{w}\geq (1-18r\varepsilon)\mathbi{e}$. If $\mathbi{w}\leq (1-18r\varepsilon)\mathbi{e}$, then using $p\lesssim n^{-1/r}$ we get
			\begin{equation}
				\label{main proof eq 3}
				n^{-\mathbi{w}/r}\cdot p^{(1-18r\varepsilon)\mathbi{e}-\mathbi{w}}\leq n^{-(1-19r\varepsilon){\bm e}/r} .
			\end{equation}
			On the other hand, if $\mathbi{w}\geq (1-18r\varepsilon)\mathbi{e}$, since we also have $\mathbi{w}\leq \mathbi{e}$, we get
			\begin{equation}
				\label{main proof eq 4}
				n^{-\mathbi{w}/r}\cdot p^{(1-18r\varepsilon)\mathbi{e}-\mathbi{w}}\leq n^{-(1-18r\varepsilon)\mathbi{e}/r}\cdot p^{-18 r\varepsilon \mathbi{e}}\leq n^{-(1-36r^{2}\varepsilon)\mathbi{e}/r},
			\end{equation}
			where in the last step we used $n^{1+1/r}p\geq 1$. 
			Combining \eqref{main proof eq 2}-\eqref{main proof eq 4} along with $\mathbi{e}\geq e_{\min}$  we obtain the second inequality in \eqref{main proof eq 1}.
			%\begin{equation}
			%\label{2nd last main}
			%      \P\left(\exists \tu{G}\subset \G(n,p):\ e(\tu{G})=\mathbi{e}\text{ and } |\W(\tu{G})|=\mathbi{w}\right)
			%     \leq \exp\left(-(1-36r^{2}\varepsilon)(1-3\varepsilon)^{1/r}\frac{\delta^{1/r}}{r}n^{1+1/r}p\log n\right).
			%\end{equation}
			
			
			Now, summing up both sides of \eqref{main proof eq 1} over allowable range of $\mathbi{e}$ and $\mathbi{w}$, as $n^{1+1/r} p \gg 1$, we derive that
			%\begin{equation}
			
			
			\begin{equation}\label{s2: bound deltanprlessthan1}
				\P\left(\bigcup_{\mathbi{e}, \mathbi{w}} \left\{\exists \tu{G}\subset \G(n,p):\ e(\tu{G})=\mathbi{e}, |\W(\tu{G})|=\mathbi{w}\right\}\right)\\[-1pt]
				%    &\leq \left(\cbar n^{1+1/r}p\log(1/p)\right)^{2}\cdot \exp\left(-(1-36r^{2}\varepsilon)(1-3\varepsilon)^{1/r}\frac{\delta^{1/r}}{r}n^{1+1/r}p\log n\right)\\[-1pt]
				\leq \exp\left(-(1-37r^{2}\varepsilon)(1-3\varepsilon)^{1/r}\frac{\delta^{1/r}}{r}n^{1+1/r}p\log n\right),
			\end{equation}
			%\end{equation}
			for all large $n$. See that this bound \eqref{s2: bound deltanprlessthan1} is optimal if $\delta np^{r}<1$. Now, assume $np^{r}\rightarrow \rho\in[1,\infty)$. In this case, we use Lemma \ref{s2: improved geb } to get that 
			$
			e(\tu{G})\geq \widehat e_{\min} \coloneqq (1-\varepsilon)\left(\left\lfloor\delta\rho(1-5\varepsilon)\right\rfloor+\left\{\delta\rho(1-5\varepsilon)\right\}^{1/r}\right)\cdot n
			$. 
			Using the lower bound $\widehat{e}_{\min}$ on the edges instead of $e_{\min}$ and repeating the same steps as in the proof of \eqref{s2: bound deltanprlessthan1} we derive \eqref{eq:p-sim-bdry}. 
			%\begin{equation*}
			%\begin{aligned}
			%    &\P\left(\bigcup_{\mathbi{e}, \mathbi{w}} \left\{\exists \tu{G}\subset \G(n,p):\ e(\tu{G})=\mathbi{e}, |\W(\tu{G})|=\mathbi{w}\right\}\right)\\[0.1pt]
			%    &\qquad \qquad \qquad \leq \exp\left(-(1-38r^{2}\varepsilon)\left(\left\lfloor\delta\rho(1-5\varepsilon)\right\rfloor+\left\{\delta\rho(1-5\varepsilon)\right\}^{1/r}\right) \frac{1}{r}n\log n\right).
			%    \end{aligned}
		%\end{equation*}
		This completes the proof.
	\end{proof}

\end{document}
